\documentclass[12pt]{article}

\usepackage[a4paper, margin=1in]{geometry}
\usepackage{amsmath,amssymb,amsthm}
\usepackage{mathtools}
\usepackage{mathrsfs}
\usepackage{enumitem}
\usepackage{xcolor}
\usepackage[hidelinks]{hyperref}
\usepackage{fontspec}

\defaultfontfeatures{Ligatures=TeX}
\IfFontExistsTF{Times New Roman}
  {\setmainfont{Times New Roman}}
  {\setmainfont{TeX Gyre Termes}}
\IfFontExistsTF{Menlo}
  {\setmonofont{Menlo}}
  {\setmonofont{Latin Modern Mono}}

\newtheorem{theorem}{Theorem}[section]
\newtheorem{proposition}[theorem]{Proposition}
\newtheorem{lemma}[theorem]{Lemma}
\newtheorem{corollary}[theorem]{Corollary}
\theoremstyle{definition}
\newtheorem{definition}[theorem]{Definition}
\theoremstyle{remark}
\newtheorem{remark}[theorem]{Remark}

\newcommand{\tightlist}{%
  \setlength{\itemsep}{0pt}\setlength{\parskip}{0pt}}

\title{%
  \textbf{Foundation III: The Perfuming Loop}\\
  \textbf{A Conditional One-Step Source-Level Feedback Program}%
}

\author{
  Sidong Liu, PhD \\
  \small iBioStratix Ltd \\
  \small \texttt{sidongliu@hotmail.com}
}

\date{July 2026}

\begin{document}
\emergencystretch=2em
\raggedbottom
\hbadness=5000
\vbadness=5000

\maketitle

\begin{abstract}
Foundation~I~\cite{FoundationI} constructs a conditional single-source
chain from a source datum \(S_{\mathrm{source}}\) to an admissible
manifest pair \((\mathcal A_c,\omega)\). Foundation~II~\cite{FoundationII}
extends this to two sources, isolating a shared manifest sector. Both
chains are one-directional: source data produce manifest sectors, but
the manifest sectors do not feed back to the source. The present paper,
Foundation~III, addresses the return arrow.

T-DOME~III~\cite{TDOMEIII} supplies a calibration loop at the ego level
under explicit residual-likelihood, convexity, drift, and hardware
assumptions: the observer's prediction-error residuals drive a
natural-gradient update of the reference-frame parameter~\(\sigma\),
with tracking controlled by a local Lyapunov estimate. Paper~III of the
Consciousness Trilogy~\cite{PaperIII} now treats the former
drift-floor theorem as withdrawn and leaves only a supplied arithmetic
comparison score for finite self-models. Neither construction closes the
loop at the source level: the ego adjusts its frame, but the decoherence
semigroup \(\{\mathcal E_t\}\) that generated the manifest sector in the
first place is left untouched.

Foundation~III defines a \emph{source-level feedback map} that lifts the
ego-level calibration signal to a controlled perturbation of the
decoherence semigroup. Under named regularity and coupling conditions,
we prove a conditional one-step closed-loop foothold with three
components. First, the feedback-modified semigroup remains a valid
source datum in the pointwise \(\sigma\)-weak Heisenberg-picture sense
used by Foundation~I v2. Second, the supplied self-description
comparison functional under the closed loop satisfies a modified
algebraic comparison bound that depends on the feedback coupling
strength. Third, the one-directional Foundation~I chain
upgrades only to a conditional one-step closed loop:
\[
\begin{aligned}
S_{\mathrm{source}}
&\longrightarrow
\mathcal O_B
\longrightarrow
\text{calibration signal}
\longrightarrow
S_{\mathrm{source}}[\text{updated}]\\
&\longrightarrow
\mathcal O_B[\text{updated}]
\longrightarrow
(\mathcal A_c,\omega)[\text{updated}].
\end{aligned}
\]
The paper does not claim that the loop eliminates the self-description
deficit. It proves only that the loop is well-posed and that the deficit
transforms in a controlled way. The question of whether iterated
perfuming converges to a fixed point, and what such a fixed point would
mean structurally, is left as an open problem.
 \end{abstract}

\clearpage
\tableofcontents
\clearpage

\section{Introduction}\label{sec:intro}

Foundation~I~\cite{FoundationI} solved the first source-side bridge in
the Consciousness Series. Starting from a source datum
\[
S_{\mathrm{source}}
=
(\mathcal H_U,\mathcal A_U,\Omega,\{\mathcal E_t\}_{t\ge 0}),
\]
it isolated a conditional chain
\[
S_{\mathrm{source}}
\to
\mathcal O_A
\to
\mathcal O_B
\to
(\mathcal A_c,\omega),
\]
where \(\mathcal O_A\) is a persistent projected record process,
\(\mathcal O_B\) is a foreground frame available under the explicit
conditional coding interface of Foundation~I, and
\((\mathcal A_c,\omega)\) is the manifest accessible pair. The chain is
one-directional: source data produce a manifest sector, and the series
then hands that sector to Paper~I of the Consciousness Trilogy
\cite{PaperI} for the geometry--modular-flow theorem and to
Paper~II~\cite{PaperII} for the higher-dimensional consistency analysis.

Foundation~II~\cite{FoundationII} extended that source-side analysis to
the first multi-source setting. There the main result was again
one-directional. Given two source chains sharing one ambient
Type~III\(_1\) substrate, one obtains a conditional inclusion theorem
for a shared admissible manifest sector. What remained absent in both
papers was any formal return arrow from manifest calibration back to the
source dynamics that generated the manifest sector in the first place.
Foundation~II enters the present paper only as programme boundary
calibration for the later multi-source feedback problem; no theorem
below imports a Foundation~II result.

The missing input now comes from two already published directions.
T-DOME~III~\cite{TDOMEIII} supplied a one-step calibration signal at the
ego level: prediction-residual statistics determine a natural-gradient
update \(\delta\sigma\), controlled by a Lyapunov tracking estimate and
penalized by the ego-rigidity term
\(\mathcal I_{\mathrm{ego}}\). Paper~III of the Consciousness Trilogy
\cite{PaperIII} proposed a local self-description floor on the manifest
side. That reconstruction-loss bound has since been withdrawn: a
lossless code for a low-dimensional orbit defeats an
\(n\)-independent positive floor. Foundation~III therefore uses its
displayed Fisher expression only as a supplied comparison functional,
not as an imported risk or reconstruction theorem. Neither that
functional nor the T-DOME calibration model, by itself, modifies the
source semigroup \(\{\mathcal E_t\}\).

The present paper addresses exactly that gap. Its aim is not to prove a
global convergence theorem for iterative self-modification or a general
iterated-feedback theorem. It does not claim that the self-description deficit
is eliminated by feedback, and it does not claim that one can derive a
feedback lift from bare multi-scale dynamics without additional
structure. Its narrower goal is to identify an honest one-step
theorem-level foothold for \emph{source-level feedback from manifest
calibration}.

The main design choice is to avoid a formal expression of the type
\[
\mathcal E_t \mapsto \mathcal E_t + \eta\,\Gamma(\delta\sigma,\mathcal E_t)
\]
as the \emph{definition} of the modified dynamics. Such a first-order
expression is useful heuristically, but it does not by itself preserve
the semigroup law or complete positivity. Instead, the paper defines a
feedback lift through an exact auxiliary semigroup
\(\{\mathscr F_r^{[\delta\sigma]}\}_{r\ge 0}\) of normal completely
positive unital maps and then sets
\[
\mathcal E_t^{[\eta]}
:=
\mathscr F_{\eta t}^{[\delta\sigma]}\circ\mathcal E_t.
\]
Once a commutation condition is imposed, the semigroup law is automatic;
the first-order feedback map \(\Gamma\) then reappears as the
\(\eta\downarrow 0\) tangent of this exact construction.

This paper therefore contributes four one-step theorem-level items.
\begin{enumerate}
\item We define a \emph{source-level feedback lift} from an admissible
  calibration signal \(\delta\sigma\) to an exact Heisenberg-picture
  semigroup perturbation on the ambient Type~III\(_1\) algebra.
\item We prove a preservation theorem: under named regularity and
  commutation inputs, the modified family
  \(\{\mathcal E_t^{[\eta]}\}_{t\ge0}\) is again a pointwise
  \(\sigma\)-weakly continuous semigroup of normal completely positive
  unital maps on \(\mathcal A_U\), so the ambient datum remains a valid
  source datum.
\item We prove a comparison theorem for a supplied local Fisher
  functional: if the chosen Foundation~I chain is feedback-stable and
  its parameters remain controlled, then that functional deforms by an
  explicit small-coupling factor \(f(\eta)\) with \(f(0)=1\). No actual
  deficit bound is inferred.
\item We package the result as a \emph{one-step return datum},
  showing that the chain
  \[
  S_{\mathrm{source}}
  \to
  \mathcal O_B
  \to
  \delta\sigma
  \to
  S_{\mathrm{source}}^{[\eta]}
  \to
  \mathcal O_B^{[\eta]}
  \to
  (\mathcal A_c^{[\eta]},\omega^{[\eta]})
  \]
  is mathematically well-posed under named inputs.
\end{enumerate}

The scope is deliberately narrow. The paper does not solve the
constructive origin of the calibration signal, which is imported from
T-DOME~III. It does not solve the full persistence problem for the
Foundation~I chain under arbitrary perturbations, which remains a
conditional interface. It does not derive explicit optimal coupling
windows or fixed points for iterated perfuming. Those are precisely the
open problems recorded later.

The paper also keeps the same Type~III discipline as Foundation~II. No
tensor-product factorization is assumed. All dynamics are written in the
Heisenberg picture on a common ambient algebra, and continuity means
pointwise \(\sigma\)-weak continuity rather than norm continuity. This
is not cosmetic: the whole point of the paper is to modify the source
semigroup without importing an external finite-dimensional controller.

The paper is organized as follows. Section~\ref{sec:feedback} defines
the feedback lift, the regularity package, and the modified source
datum. Section~\ref{sec:preservation} proves the source-datum
preservation theorem. Section~\ref{sec:deficit} states the extra
comparison inputs needed to pass through Foundation~I and then proves
the comparison-functional theorem. Section~\ref{sec:closed} packages
the resulting one-step return map. Section~\ref{sec:open} records the
open problems. Section~\ref{sec:ledger} collects the named conditional
inputs. Appendix~\ref{app:interfaces} restates the conditional
Foundation~I and T-DOME~III interfaces and records the withdrawal of the
former Paper~III floor.

\section{Feedback Lift and Regularity Conditions}\label{sec:feedback}

The first task is to isolate the additional object that Foundation~III
needs and the exact conditions under which it is allowed to act on the
source semigroup.

\begin{definition}[Admissible calibration signal]\label{def:signal}
Fix a source datum \(S_{\mathrm{source}}\) together with a chosen
Foundation~I chain up to Stage~B, so that one has
\[
\mathcal O_B
=
(\mathcal P,\omega_{\mathcal P,t},\mathcal K_*(t,s),\mathfrak F,\pi_{\mathrm{ego}}).
\]
An \emph{admissible calibration signal} is a vector
\[
\delta\sigma \in T_\sigma\Sigma
\]
in the local calibration parameter space of T-DOME~III such that:
\begin{enumerate}
\item \(\delta\sigma\) is obtained from the natural-gradient calibration
  step imported from T-DOME~III;
\item \(\sigma+\delta\sigma\) remains inside the Lyapunov neighborhood on
  which the T-DOME~III tracking estimate is valid;
\item there exists a bound
  \(\|\delta\sigma\| \leq R_{\mathrm{cal}}\) for some finite
  calibration radius \(R_{\mathrm{cal}}>0\).
\end{enumerate}
\end{definition}

\begin{remark}[One-step rather than iterative calibration]
Definition~\ref{def:signal} is intentionally local. The paper treats a
single feedback step driven by one calibration signal
\(\delta\sigma\). An iterated sequence
\(\delta\sigma_0,\delta\sigma_1,\dots\) would require a separate fixed
point or long-time theory, and nothing in the present paper should be
read as supplying that theory.
\end{remark}

\begin{definition}[Feedback regularity package]\label{def:FB}
Let \(S_{\mathrm{source}}\) be a source datum and let
\(\delta\sigma\) be an admissible calibration signal. A
\emph{feedback regularity package} consists of a family
\[
\mathscr F^{[\delta\sigma]}
=
\{\mathscr F_r^{[\delta\sigma]}\}_{r\ge0}
\]
of normal maps on \(\mathcal A_U\) satisfying the following conditions.
\begin{enumerate}
\item \textbf{FB-A: CP source lift.}
  \(\{\mathscr F_r^{[\delta\sigma]}\}_{r\ge0}\) is a pointwise
  \(\sigma\)-weakly continuous semigroup of normal completely positive
  unital maps on \(\mathcal A_U\), with
  \[
  \mathscr F_0^{[\delta\sigma]}=\mathrm{id}_{\mathcal A_U},
  \qquad
  \mathscr F_r^{[0]}=\mathrm{id}_{\mathcal A_U}
  \quad
  \forall r\ge0.
  \]
\item \textbf{FB-B: small-signal differentiability.}
  For every \(A\in\mathcal A_U\), the map
  \(r\mapsto \mathscr F_r^{[\delta\sigma]}(A)\) is
  \(\sigma\)-weakly differentiable at \(r=0\). We denote the
  derivative by
  \[
  \mathfrak G_{\delta\sigma}(A)
  :=
  \left.
  \frac{d}{dr}\mathscr F_r^{[\delta\sigma]}(A)
  \right|_{r=0},
  \]
  and record the normalization
  \(\mathfrak G_{\delta\sigma}(\mathbf 1)=0\), which follows from the
  unitality in FB-A because
  \(\mathscr F_r^{[\delta\sigma]}(\mathbf 1)=\mathbf 1\) for all
  \(r\ge0\).
\item \textbf{FB-C: chain compatibility.}
  The feedback lift commutes with the ambient source semigroup:
  \[
  \mathscr F_r^{[\delta\sigma]}\mathcal E_t
  =
  \mathcal E_t\mathscr F_r^{[\delta\sigma]}
  \qquad
  \forall r,t\ge0.
  \]
  In addition, for the chosen Foundation~I chain, the maps
  \(\mathscr F_r^{[\delta\sigma]}\) preserve the named structural inputs
  that are required later to reconstruct Stage~A, Stage~B, and the
  manifest pair.
\end{enumerate}
\end{definition}

\begin{remark}[Why the feedback acts through an exact semigroup]
The exact feedback semigroup in Definition~\ref{def:FB} is not an
ornament. A first-order perturbation
\(\mathcal E_t+\eta\,\Gamma(\delta\sigma,\mathcal E_t)\) does not
automatically preserve the semigroup property or complete positivity.
The exact semigroup lift is the cleanest way to keep those properties in
the Heisenberg-picture Type~III setting. The first-order map
\(\Gamma\) reappears only after the exact object has been defined.
\end{remark}

\begin{remark}[Meaning of the strict commutation input]
The commutation requirement in FB-C is strong and should be read
carefully. It does \emph{not} say that feedback is dynamically trivial.
For \(\eta>0\),
\[
\mathcal E_t^{[\eta]}
=
\mathscr F_{\eta t}^{[\delta\sigma]}\circ \mathcal E_t
\]
generally differs from \(\mathcal E_t\) unless the calibration signal is
zero or the coupling is turned off. The role of strict commutation is
only to keep the algebra-level time ordering unchanged, so that the
modified dynamics remains an exact semigroup without entering a more
delicate generator-level perturbation problem. Relaxing FB-C to an
approximate-commutation condition is left open.
\end{remark}

\begin{proposition}[A nontrivial commuting-dephasing feedback lift]
\label{prop:fdlift}
The package FB-A--FB-C has a finite-dimensional normal UCP realization
that depends nontrivially on \(\delta\sigma\) and changes the source
dynamics beyond a relabeling.
\end{proposition}

\begin{proof}
Take \(\mathcal A_U=M_2\otimes M_2\), let \(Z_1,Z_2\) be Pauli \(Z\) on
the two factors, and write \(\operatorname{Ad}_{Z_j}(A)=Z_jAZ_j\). Fix
\(\gamma,\kappa>0\) and a continuous scalar response
\[
q(\delta\sigma):=\|\delta\sigma\|^2.
\]
Define bounded generators
\[
\mathcal L_0=\gamma(\operatorname{Ad}_{Z_1}-\mathrm{id}),
\qquad
\mathfrak G_{\delta\sigma}
=\kappa q(\delta\sigma)
(\operatorname{Ad}_{Z_2}-\mathrm{id}),
\]
and semigroups
\[
\mathcal E_t=e^{t\mathcal L_0},
\qquad
\mathscr F_r^{[\delta\sigma]}=e^{r\mathfrak G_{\delta\sigma}}.
\]
More explicitly,
\[
\mathscr F_r^{[\delta\sigma]}(A)
=a_r A+(1-a_r)Z_2AZ_2,
\qquad
a_r=\frac{1+e^{-2\kappa q(\delta\sigma)r}}{2}.
\]
Thus \(\mathscr F_r^{[\delta\sigma]}\) is a normal UCP unital
semigroup, equals the identity when \(\delta\sigma=0\), and has the
stated generator at \(r=0\). The same formula with
\((Z_2,\kappa q,r)\) replaced by \((Z_1,\gamma,t)\) verifies these
properties for \(\mathcal E_t\).

Because \(\operatorname{Ad}_{Z_1}\) and
\(\operatorname{Ad}_{Z_2}\) commute, so do the two semigroups, which is
the algebraic clause of FB-C. A Foundation~I input package supported on
the joint diagonal algebra is fixed by both dephasings, so the named
structural inputs can be chosen feedback-stable in this model. Yet for
\(A=\mathbf1\otimes X\),
\[
\mathcal E_t(A)=A,
\qquad
\mathcal E_t^{[\eta]}(A)
=e^{-2\eta\kappa q(\delta\sigma)t}A.
\]
For \(\eta>0\) and \(\delta\sigma\ne0\), the lift therefore adds a real
decay channel that the source semigroup did not have; it is not an inert
renaming. This is an FD/I witness of compatibility, not a construction
of a Type~III feedback lift.
\end{proof}

\begin{remark}[What strict commutation can and cannot add]
Proposition~\ref{prop:fdlift} shows that FB-C does not force the lift to
be the identity. It also displays the limitation: the commuting
generators are simultaneously diagonal on joint Pauli sectors, so
feedback only reweights pre-existing joint modes and does not mix them.
Creating new couplings between source sectors generally requires a
near-commuting or time-ordered theory, which remains open.
\end{remark}

\begin{definition}[Small-coupling regime and modified semigroup]
\label{def:modified}
Fix a feedback regularity package
\(\mathscr F^{[\delta\sigma]}\). A coupling strength \(\eta\ge0\) is in
the \emph{small-coupling regime} if \(0\le \eta < \eta_*\) for some
window \(\eta_*>0\) on which the later comparison inputs remain valid.
For a first-order interpretation on a return horizon \([0,T]\), we
additionally require \(\eta T\ll1\). The exact composed semigroup below
exists without this perturbative estimate, but its tangent need not be
uniformly small at long times.
For such \(\eta\), define the \emph{feedback-modified semigroup} by
\[
\mathcal E_t^{[\eta]}
:=
\mathscr F_{\eta t}^{[\delta\sigma]}\circ \mathcal E_t
\qquad
(t\ge0).
\]
The corresponding \emph{modified source datum} is
\[
S_{\mathrm{source}}^{[\eta]}
:=
(\mathcal H_U,\mathcal A_U,\Omega,\{\mathcal E_t^{[\eta]}\}_{t\ge0}).
\]
\end{definition}

\begin{remark}[Heisenberg-picture convention]
As in Foundation~II, all semigroups in Definition~\ref{def:modified}
act on observables. Pointwise continuity therefore means
\(\sigma\)-weak continuity on \(\mathcal A_U\), not norm continuity.
Nothing in the present paper uses or claims a norm-continuous
generator-level perturbation theorem.
\end{remark}

\begin{definition}[First-order feedback map]\label{def:Gamma}
Under FB-B, define the \emph{first-order feedback map} by
\[
\Gamma(\delta\sigma,\mathcal E_t)(A)
:=
t\,\mathfrak G_{\delta\sigma}(\mathcal E_t(A)),
\qquad
A\in\mathcal A_U.
\]
Equivalently, for each fixed \(t\ge0\),
\[
\mathcal E_t^{[\eta]}(A)
=
\mathcal E_t(A)
+
\eta\,\Gamma(\delta\sigma,\mathcal E_t)(A)
+
o(\eta)
\]
in the \(\sigma\)-weak topology as \(\eta\downarrow 0\).
\end{definition}

\begin{proof}[Verification of the expansion]
By Definition~\ref{def:modified},
\[
\mathcal E_t^{[\eta]}(A)
=
\mathscr F_{\eta t}^{[\delta\sigma]}(\mathcal E_t(A)).
\]
Applying the \(r=0\) derivative from FB-B at \(r=\eta t\) gives
\[
\mathscr F_{\eta t}^{[\delta\sigma]}(\mathcal E_t(A))
=
\mathcal E_t(A)
+
\eta t\,\mathfrak G_{\delta\sigma}(\mathcal E_t(A))
+
o(\eta),
\]
which is exactly the stated formula.
\end{proof}

\begin{remark}[What Definition~\ref{def:Gamma} does not prove]
The first-order map \(\Gamma\) is merely the tangent of the exact
feedback lift. Definition~\ref{def:Gamma} does not claim that an
arbitrary bounded normal map \(\Gamma\) integrates to a completely
positive semigroup, nor does it characterize the most general
generator-level perturbations compatible with Type~III dynamics. Those
questions are left open.
\end{remark}

\begin{remark}[The tangent is not uniform at long times]
The identity
\[
\Gamma(\delta\sigma,\mathcal E_t)(A)
=t\,\mathfrak G_{\delta\sigma}(\mathcal E_t(A))
\]
contains an explicit factor of \(t\). Thus the first-order perturbation
is controlled for fixed \(t\) as \(\eta\downarrow0\), or on windows with
\(\eta t\ll1\) under a uniform generator bound. Small \(\eta\) alone
does not give a perturbation uniformly small for arbitrarily long times.
This is one reason the present result does not iterate to a long-time
perfuming theorem.
\end{remark}

\section{Source-Datum Preservation}\label{sec:preservation}

We now prove the first of the paper's three main results.

\begin{lemma}[Semigroup law]\label{lem:semigroup}
Assume FB-A--FB-C. Then the family
\(\{\mathcal E_t^{[\eta]}\}_{t\ge0}\) from
Definition~\ref{def:modified} satisfies
\[
\mathcal E_t^{[\eta]}\mathcal E_s^{[\eta]}
=
\mathcal E_{t+s}^{[\eta]}
\qquad
\forall s,t\ge0.
\]
\end{lemma}

\begin{proof}
For \(A\in\mathcal A_U\),
\begin{align*}
\mathcal E_t^{[\eta]}\mathcal E_s^{[\eta]}(A)
&=
\mathscr F_{\eta t}^{[\delta\sigma]}\,
\mathcal E_t
\Bigl(
\mathscr F_{\eta s}^{[\delta\sigma]}\,
\mathcal E_s(A)
\Bigr) \\
&=
\mathscr F_{\eta t}^{[\delta\sigma]}
\mathscr F_{\eta s}^{[\delta\sigma]}
\mathcal E_t\mathcal E_s(A),
\end{align*}
where FB-C is used to commute
\(\mathcal E_t\) past \(\mathscr F_{\eta s}^{[\delta\sigma]}\). Since
\(\{\mathscr F_r^{[\delta\sigma]}\}\) and \(\{\mathcal E_t\}\) are both
semigroups,
\[
\mathscr F_{\eta t}^{[\delta\sigma]}
\mathscr F_{\eta s}^{[\delta\sigma]}
\mathcal E_t\mathcal E_s
=
\mathscr F_{\eta(t+s)}^{[\delta\sigma]}
\mathcal E_{t+s}
=
\mathcal E_{t+s}^{[\eta]}.
\]
\end{proof}

\begin{lemma}[Pointwise \(\sigma\)-weak continuity and positivity]
\label{lem:continuity}
Assume FB-A--FB-C. Then for every \(A\in\mathcal A_U\),
\[
t\longmapsto \mathcal E_t^{[\eta]}(A)
\]
is pointwise \(\sigma\)-weakly continuous. Moreover, each
\(\mathcal E_t^{[\eta]}\) is normal, completely positive, and unital.
\end{lemma}

\begin{proof}
Fix \(A\in\mathcal A_U\) and a normal functional
\(\varphi\in(\mathcal A_U)_*\). By Definition~\ref{def:modified},
\[
\varphi(\mathcal E_t^{[\eta]}(A))
=
\varphi\!\left(
\mathscr F_{\eta t}^{[\delta\sigma]}(\mathcal E_t(A))
\right).
\]
The map \(t\mapsto \mathcal E_t(A)\) is \(\sigma\)-weakly continuous by
the source-datum assumption, and
\(r\mapsto \mathscr F_r^{[\delta\sigma]}(B)\) is
\(\sigma\)-weakly continuous for every fixed \(B\in\mathcal A_U\) by
FB-A. To prove continuity of the composition, fix \(t_0\ge0\) and write
\begin{align*}
&\left|
\varphi\bigl(\mathscr F_{\eta t}^{[\delta\sigma]}(\mathcal E_t(A))\bigr)
-
\varphi\bigl(\mathscr F_{\eta t_0}^{[\delta\sigma]}(\mathcal E_{t_0}(A))\bigr)
\right| \\
&\qquad \le
\left|
\varphi\bigl(\mathscr F_{\eta t}^{[\delta\sigma]}(\mathcal E_t(A))\bigr)
-
\varphi\bigl(\mathscr F_{\eta t}^{[\delta\sigma]}(\mathcal E_{t_0}(A))\bigr)
\right| \\
&\qquad\quad +
\left|
\varphi\bigl(\mathscr F_{\eta t}^{[\delta\sigma]}(\mathcal E_{t_0}(A))\bigr)
-
\varphi\bigl(\mathscr F_{\eta t_0}^{[\delta\sigma]}(\mathcal E_{t_0}(A))\bigr)
\right|.
\end{align*}
For the second term, \(\mathcal E_{t_0}(A)\) is fixed, so pointwise
\(\sigma\)-weak continuity of
\(\{\mathscr F_r^{[\delta\sigma]}\}_{r\ge0}\) gives convergence to
zero as \(t\to t_0\).

For the first term, use the predual semigroup
\(\{\mathscr F_{r,*}^{[\delta\sigma]}\}_{r\ge0}\). Since
\(\{\mathscr F_r^{[\delta\sigma]}\}\) is a pointwise
\(\sigma\)-weakly continuous semigroup of normal maps, the predual
family is norm-continuous on \((\mathcal A_U)_*\). Thus
\begin{align*}
&\left|
\varphi\bigl(\mathscr F_{\eta t}^{[\delta\sigma]}(\mathcal E_t(A))\bigr)
-
\varphi\bigl(\mathscr F_{\eta t}^{[\delta\sigma]}(\mathcal E_{t_0}(A))\bigr)
\right| \\
&\qquad =
\left|
\mathscr F_{\eta t,*}^{[\delta\sigma]}(\varphi)
\bigl(\mathcal E_t(A)-\mathcal E_{t_0}(A)\bigr)
\right| \\
&\qquad \le
\left|
\bigl(\mathscr F_{\eta t,*}^{[\delta\sigma]}(\varphi)
-
\mathscr F_{\eta t_0,*}^{[\delta\sigma]}(\varphi)\bigr)
\bigl(\mathcal E_t(A)-\mathcal E_{t_0}(A)\bigr)
\right| \\
&\qquad\quad +
\left|
\mathscr F_{\eta t_0,*}^{[\delta\sigma]}(\varphi)
\bigl(\mathcal E_t(A)-\mathcal E_{t_0}(A)\bigr)
\right|.
\end{align*}
The first summand tends to zero because
\(\mathscr F_{\eta t,*}^{[\delta\sigma]}(\varphi)\to
\mathscr F_{\eta t_0,*}^{[\delta\sigma]}(\varphi)\) in predual norm,
while \(\mathcal E_t(A)-\mathcal E_{t_0}(A)\) is norm-bounded by
\(2\|A\|\). The second summand tends to zero because
\(\mathcal E_t(A)\to\mathcal E_{t_0}(A)\) \(\sigma\)-weakly and
\(\mathscr F_{\eta t_0,*}^{[\delta\sigma]}(\varphi)\) is a fixed normal
functional. Hence the first term also tends to zero.

Hence \(\{\mathcal E_t^{[\eta]}\}\) is pointwise \(\sigma\)-weakly
continuous.

Normality and complete positivity are preserved under composition of
normal completely positive maps, and unitality is preserved under
composition of unital maps. Since both
\(\mathscr F_{\eta t}^{[\delta\sigma]}\) and \(\mathcal E_t\) have
those properties, so does \(\mathcal E_t^{[\eta]}\).
\end{proof}

\begin{theorem}[Source-datum preservation]\label{thm:preserve}
Let
\[
S_{\mathrm{source}}
=
(\mathcal H_U,\mathcal A_U,\Omega,\{\mathcal E_t\}_{t\ge0})
\]
be a source datum in the sense imported from Foundation~I, let
\(\delta\sigma\) be an admissible calibration signal, and assume that a
feedback regularity package FB-A--FB-C is given. Then for every
\(\eta\in[0,\eta_*)\), the modified datum
\[
S_{\mathrm{source}}^{[\eta]}
=
(\mathcal H_U,\mathcal A_U,\Omega,\{\mathcal E_t^{[\eta]}\}_{t\ge0})
\]
is again a valid source datum: the family
\(\{\mathcal E_t^{[\eta]}\}_{t\ge0}\) is a pointwise
\(\sigma\)-weakly continuous semigroup of normal completely positive
unital maps on \(\mathcal A_U\), and \(\Omega\) remains cyclic and
separating for \(\mathcal A_U\).
\end{theorem}

\begin{proof}
Lemma~\ref{lem:semigroup} gives the semigroup law and
Lemma~\ref{lem:continuity} gives pointwise \(\sigma\)-weak continuity,
normality, complete positivity, and unitality.

The ambient standard-form triple
\((\mathcal H_U,\mathcal A_U,\Omega)\) is not changed by the feedback
step. In particular, \(\Omega\) remains cyclic and separating for the
same von Neumann algebra \(\mathcal A_U\). Therefore
\(S_{\mathrm{source}}^{[\eta]}\) satisfies the same source-datum format
as \(S_{\mathrm{source}}\).
\end{proof}

\begin{remark}[What Theorem~\ref{thm:preserve} does not prove]
The theorem proves only that the feedback-modified dynamics remains in
the theorem domain of the Foundation~I source-datum definition. It does
not prove that the modified datum automatically reproduces the same
persistent projection, the same foreground frame, or the same manifest
algebra. Those are downstream interface questions and require additional
inputs.
\end{remark}

\begin{corollary}[Exact CP lift with first-order tangent]
\label{cor:firstorder}
Under the hypotheses of Theorem~\ref{thm:preserve}, the map
\(\Gamma(\delta\sigma,\mathcal E_t)\) from
Definition~\ref{def:Gamma} is the \(\eta\downarrow 0\) tangent of an
exact Heisenberg-picture completely positive perturbation.
\end{corollary}

\begin{proof}
This is immediate from Theorem~\ref{thm:preserve} together with the
first-order expansion of Definition~\ref{def:Gamma}.
\end{proof}

\section{Modified Local Comparison Functional}\label{sec:deficit}

The preservation theorem keeps the modified dynamics inside the
source-datum domain, but it does not yet produce a manifest pair. For
that we need the chosen Foundation~I chain to survive the perturbation
and the supplied local comparison parameters to remain controlled.

\begin{definition}[Deficit-comparison package]\label{def:DFA}
Fix a source datum \(S_{\mathrm{source}}\), a chosen Foundation~I chain
from \(S_{\mathrm{source}}\) to \((\mathcal A_c,\omega)\), and a
feedback-modified datum \(S_{\mathrm{source}}^{[\eta]}\). We say that
\emph{DF-A} holds on the coupling window \([0,\eta_*)\) if:
\begin{enumerate}
\item \textbf{DF-A(1): feedback-stable Foundation~I chain.}
  For every \(\eta\in[0,\eta_*)\), the named conditional inputs used in
  the chosen Foundation~I chain remain available for
  \(S_{\mathrm{source}}^{[\eta]}\). Hence one obtains a modified manifest
  pair
  \[
  (\mathcal A_c^{[\eta]},\omega^{[\eta]}).
  \]
\item \textbf{DF-A(2): Fisher control.}
  There exist continuous functions
  \(\alpha_F,\beta_{\mathrm{ego}}:[0,\eta_*)\to(0,\infty)\) with
  \[
  \alpha_F(0)=\beta_{\mathrm{ego}}(0)=1
  \]
  such that
  \[
  \mathcal I_F^{[\eta]}
  \ge
  \alpha_F(\eta)\,\mathcal I_F,
  \qquad
  \mathcal I_{\mathrm{ego}}^{[\eta]}
  \le
  \beta_{\mathrm{ego}}(\eta)\,\mathcal I_{\mathrm{ego}},
  \]
  where \(\mathcal I_F\) and
  \(\mathcal I_{\mathrm{ego}}\) are the baseline manifest-side Fisher
  sensitivity and ego-rigidity constants, and
  \(\mathcal I_F^{[\eta]}\),
  \(\mathcal I_{\mathrm{ego}}^{[\eta]}\) are the corresponding modified
  quantities. Assume also
  \[
  \mathcal I_F+\mathcal I_{\mathrm{ego}}>0,
  \qquad
  \mathcal I_F^{[\eta]}+\mathcal I_{\mathrm{ego}}^{[\eta]}>0,
  \]
  so the comparison functionals are defined.
\item \textbf{DF-A(3): supplied local comparison parameters.}
  A controlled neighborhood is equipped with a stipulated parameter
  \(\rho(\eta)\in[0,\tfrac12)\), continuous in \(\eta\), with baseline
  value \(\rho(0)=:\rho_0\). This parameter and the Fisher quantities
  define the comparison functional below; they do not import a valid
  Paper~III reconstruction-loss theorem.
\end{enumerate}
\end{definition}

\begin{proposition}[Passage from the modified source datum to the modified manifest pair]
\label{prop:chain}
Assume the hypotheses of Theorem~\ref{thm:preserve} and DF-A(1). Then
for every \(\eta\in[0,\eta_*)\), the modified source datum
\(S_{\mathrm{source}}^{[\eta]}\) admits a conditional Foundation~I chain
\[
S_{\mathrm{source}}^{[\eta]}
\to
\mathcal O_A^{[\eta]}
\to
\mathcal O_B^{[\eta]}
\to
(\mathcal A_c^{[\eta]},\omega^{[\eta]}).
\]
\end{proposition}

\begin{proof}
Theorem~\ref{thm:preserve} puts
\(S_{\mathrm{source}}^{[\eta]}\) inside the theorem domain of the
Foundation~I source-datum definition. DF-A(1) states exactly that the
named conditional inputs used by the chosen Foundation~I chain remain
available after the feedback perturbation. The Foundation~I conditional
theorem can therefore be re-run with the modified datum, yielding the
stated outputs.
\end{proof}

\begin{definition}[Baseline and modified local comparison functionals]
\label{def:floor}
Under DF-A, define
\[
\mathfrak b_{\mathrm{loc}}
:=
\frac{(1-2\rho_0)\,\mathcal I_F}
{2\bigl(\mathcal I_F+\mathcal I_{\mathrm{ego}}\bigr)},
\]
and
\[
\mathfrak b_{\mathrm{loc}}^{[\eta]}
:=
\frac{(1-2\rho(\eta))\,\mathcal I_F^{[\eta]}}
{2\bigl(\mathcal I_F^{[\eta]}+\mathcal I_{\mathrm{ego}}^{[\eta]}\bigr)}.
\]
These are algebraic comparison scores. The formula is retained from the
former Paper~III interface, but neither score is asserted here to lower
bound reconstruction loss, estimation risk, or a minimax deficit.
\end{definition}

\begin{theorem}[Comparison of supplied local functionals]\label{thm:deficit}
Assume the hypotheses of Proposition~\ref{prop:chain} and the full
package DF-A(1)--DF-A(3). Then for every \(\eta\in[0,\eta_*)\),
\[
\mathfrak b_{\mathrm{loc}}^{[\eta]}
\ge
f(\eta)\,
\mathfrak b_{\mathrm{loc}},
\]
where
\[
f(\eta)
:=
\frac{1-2\rho(\eta)}{1-2\rho_0}\,
\frac{\alpha_F(\eta)
\bigl(\mathcal I_F+\mathcal I_{\mathrm{ego}}\bigr)}
{\alpha_F(\eta)\mathcal I_F
+\beta_{\mathrm{ego}}(\eta)\mathcal I_{\mathrm{ego}}}.
\]
In particular, \(f(0)=1\), and by continuity \(f(\eta)>0\) for all
sufficiently small \(\eta\).
\end{theorem}

\begin{proof}
By Proposition~\ref{prop:chain}, the modified datum yields a manifest
pair \((\mathcal A_c^{[\eta]},\omega^{[\eta]})\). For fixed
\(c>0\) and \(d\ge0\), the function \(x\mapsto cx/(x+d)\) is increasing
where the denominator is nonzero. Therefore
substituting the lower bound for \(\mathcal I_F^{[\eta]}\) lowers the
whole fraction, rather than creating a denominator sign reversal.
Likewise the fraction is decreasing in its nonnegative second argument.
Using DF-A(2),
\[
\mathcal I_F^{[\eta]}
\ge
\alpha_F(\eta)\mathcal I_F,
\qquad
\mathcal I_{\mathrm{ego}}^{[\eta]}
\le
\beta_{\mathrm{ego}}(\eta)\mathcal I_{\mathrm{ego}},
\]
so
\[
\mathfrak b_{\mathrm{loc}}^{[\eta]}
\ge
\frac{(1-2\rho(\eta))\,\alpha_F(\eta)\mathcal I_F}
{2\bigl(\alpha_F(\eta)\mathcal I_F
+\beta_{\mathrm{ego}}(\eta)\mathcal I_{\mathrm{ego}}\bigr)}.
\]
Now factor out the baseline score from
Definition~\ref{def:floor}:
\begin{align*}
\mathfrak b_{\mathrm{loc}}^{[\eta]}
&\ge
\frac{1-2\rho(\eta)}{1-2\rho_0}\,
\frac{\alpha_F(\eta)
\bigl(\mathcal I_F+\mathcal I_{\mathrm{ego}}\bigr)}
{\alpha_F(\eta)\mathcal I_F
+\beta_{\mathrm{ego}}(\eta)\mathcal I_{\mathrm{ego}}}
\cdot
\frac{(1-2\rho_0)\mathcal I_F}
{2(\mathcal I_F+\mathcal I_{\mathrm{ego}})} \\
&=
f(\eta)\,
\mathfrak b_{\mathrm{loc}}.
\end{align*}
The formula for \(f(0)\) follows from
\(\rho(0)=\rho_0\) and \(\alpha_F(0)=\beta_{\mathrm{ego}}(0)=1\).
Continuity then gives positivity for sufficiently small \(\eta\).
\end{proof}

\begin{remark}[No monotonicity claim for \(f(\eta)\)]
The theorem asserts only that \(f(0)=1\) and that \(f(\eta)\) remains
positive for sufficiently small coupling. No monotonicity statement is
claimed. Depending on the detailed response of the manifest Fisher
sensitivity and the ego-rigidity term, one may have \(f(\eta)<1\) or
\(f(\eta)>1\). In particular, a feedback step that increases
\(\mathcal I_F^{[\eta]}\) or decreases
\(\mathcal I_{\mathrm{ego}}^{[\eta]}\) can raise the score rather than
lower it.
\end{remark}

\begin{remark}[The former actual-deficit corollary is withdrawn]
No saturation assumption can repair the missing premise that
\(\mathfrak b_{\mathrm{loc}}\) is a valid positive lower bound on the
actual reconstruction loss. Establishing such a bridge requires either
a specified sampling/measurement model and an estimation-risk theorem,
or a genuine rate-distortion bottleneck. The present paper supplies
neither, so it makes no inequality claim for an actual minimax deficit.
\end{remark}

\begin{remark}[What Theorem~\ref{thm:deficit} does not prove]
The theorem gives a controlled deformation of a supplied algebraic
score. It does not prove that the true minimax deficit is monotone in
\(\eta\), it does not prove that feedback reduces the deficit, and it
does not prove that the deficit can be driven to zero. Those stronger
claims would require a valid risk/distortion model and a transport theory
for optimal self-models.
\end{remark}

\begin{remark}[Why the comparison factor is the right theorem target]
The factor \(f(\eta)\) isolates exactly what feedback changes at theorem
level:
\begin{enumerate}
\item the supplied local coefficient \(\rho(\eta)\);
\item the manifest Fisher sensitivity
  \(\mathcal I_F^{[\eta]}\);
\item the ego-rigidity contribution
  \(\mathcal I_{\mathrm{ego}}^{[\eta]}\).
\end{enumerate}
No stronger conclusion is claimed. In particular, the theorem does not
say that a better calibration signal necessarily increases
\(\mathcal I_F^{[\eta]}\) or decreases
\(\mathcal I_{\mathrm{ego}}^{[\eta]}\); it only says that if those
quantities remain controlled, then the comparison score remains
controlled as well.
\end{remark}

\section{Conditionally Well-Posed One-Step Return Map}\label{sec:closed}

We can now state the paper's final theorem.

\begin{definition}[One-step return datum]\label{def:closed}
Under the hypotheses of Theorem~\ref{thm:deficit}, the
\emph{one-step return datum at coupling \(\eta\)} is the tuple
\[
\mathfrak C_\eta
:=
\bigl(
S_{\mathrm{source}},
\mathcal O_B,
(\mathcal A_c,\omega),
\delta\sigma,
S_{\mathrm{source}}^{[\eta]},
\mathcal O_B^{[\eta]},
(\mathcal A_c^{[\eta]},\omega^{[\eta]})
\bigr),
\]
where \(\mathcal O_B\) and \((\mathcal A_c,\omega)\) are the baseline
Stage~B and Stage~C outputs of the chosen Foundation~I chain,
\(\delta\sigma\) is the admissible calibration signal extracted from the
retained Stage~B data, and
\(\mathcal O_B^{[\eta]}\) together with
\((\mathcal A_c^{[\eta]},\omega^{[\eta]})\) are the corresponding
modified outputs after the feedback step.
\end{definition}

\begin{theorem}[Conditional one-step return-map well-posedness]\label{thm:closed}
Let \(S_{\mathrm{source}}\) be a source datum, let
\(\delta\sigma\) be an admissible calibration signal, and assume the
packages FB-A--FB-C and DF-A(1)--DF-A(3). Then for every sufficiently
small \(\eta\in[0,\eta_*)\), and on any fixed return horizon \([0,T]\)
with \(\eta T\ll1\), the one-step return datum
\(\mathfrak C_\eta\) of Definition~\ref{def:closed} is well-defined.
Equivalently, the diagram
\[
S_{\mathrm{source}}
\longrightarrow
 \mathcal O_B
\longrightarrow
\delta\sigma
\longrightarrow
S_{\mathrm{source}}^{[\eta]}
\longrightarrow
 \mathcal O_B^{[\eta]}
 \longrightarrow
 (\mathcal A_c^{[\eta]},\omega^{[\eta]})
\]
is a conditionally well-posed one-step return into the source-datum
category.
\end{theorem}

\begin{proof}
The baseline arrow
\[
S_{\mathrm{source}}\to \mathcal O_B \to (\mathcal A_c,\omega)
\]
is the chosen Foundation~I chain. The arrow
\[
\mathcal O_B \to\delta\sigma
\]
is the imported T-DOME~III calibration interface. By
Theorem~\ref{thm:preserve}, the feedback step
\[
\delta\sigma \to S_{\mathrm{source}}^{[\eta]}
\]
is well-defined at the source-datum level. By
Proposition~\ref{prop:chain}, the same conditional Foundation~I chain
can then be run again on \(S_{\mathrm{source}}^{[\eta]}\), yielding
\[
S_{\mathrm{source}}^{[\eta]}
\to
\mathcal O_B^{[\eta]}
\to
(\mathcal A_c^{[\eta]},\omega^{[\eta]}).
\]
Thus every component of \(\mathfrak C_\eta\) is defined. The supplied
comparison score is controlled by Theorem~\ref{thm:deficit}. This does
not turn the one-step return into a self-consistent loop or an actual
deficit theorem.
\end{proof}

\begin{corollary}[Packaged one-step return map]
\label{cor:map}
Under the hypotheses of Theorem~\ref{thm:closed}, there is a well-defined
map
\[
\mathscr L_\eta:
\bigl(S_{\mathrm{source}},\delta\sigma\bigr)
\longmapsto
\bigl(
S_{\mathrm{source}}^{[\eta]},
\mathcal O_B^{[\eta]},
(\mathcal A_c^{[\eta]},\omega^{[\eta]}),
\mathfrak b_{\mathrm{loc}}^{[\eta]}
\bigr).
\]
\end{corollary}

\begin{proof}
This is just a repackaging of Theorem~\ref{thm:closed} and
Theorem~\ref{thm:deficit}.
\end{proof}

\begin{remark}[What Theorem~\ref{thm:closed} does not prove]
The theorem is a one-step well-posedness statement only. It does not
claim that repeated application of \(\mathscr L_\eta\) converges, that a
fixed point exists, that a fixed point would be unique, or that the
return map is thermodynamically optimal. It also does not address the
multi-source feedback problem in which several manifest sectors update a
shared ambient semigroup simultaneously.
\end{remark}

\section{Open Problems}\label{sec:open}

The present note defines only a conditionally well-posed one-step return.
The actual programme still leaves several substantive gaps.

\begin{enumerate}
\item \textbf{Constructive origin of a Type~III feedback lift.}
  Proposition~\ref{prop:fdlift} supplies a nontrivial FD/I witness, but
  the general theorem still takes the source-level semigroup lift
  \(\mathscr F^{[\delta\sigma]}\) as named input. A future result must
  derive a normal Type~III lift from the calibration signal itself.
\item \textbf{Generator-level perturbation theory.} The present paper
  works with an exact auxiliary semigroup because it preserves complete
  positivity and the semigroup law cleanly. A sharper generator-level
  perturbation theorem on \(W^*\)-algebras, closer to the existing
  structural theory of dynamical semigroups on \(W^*\)-algebras
  \cite{FrigerioVerri1982}, would have to control densely defined
  predual generators, domain stability, closability, and resolvent
  bounds in a regime where pointwise \(\sigma\)-weak continuity is the
  natural topology and norm continuity is usually unavailable. Naive
  first-order generator perturbations may fail to preserve complete
  positivity, unitality, or even semigroup well-posedness.
\item \textbf{Weakening the commutation input.} FB-C requires strict
  algebra-level commutation between the feedback lift and the source
  semigroup. A weaker theorem based on approximate commutation,
  asymptotic commutation, or commutation only on the relevant stable
  subalgebras would be more flexible and physically more natural, but it
  would need a replacement for the exact time-ordering argument used in
  Section~\ref{sec:preservation}. The finite witness shows explicitly
  that strict commutation can change decay rates while still only
  reweighting joint spectral sectors.
\item \textbf{Quantitative coupling thresholds.} The small-coupling
  window \([0,\eta_*)\) is left abstract. The programme still lacks an
  explicit estimate of \(\eta_*\) in terms of Lyapunov constants,
  semigroup norms, and the local BKM/Fisher geometry.
\item \textbf{Persistence of the Foundation~I chain.} DF-A(1) is a real
  interface burden. One needs a theorem stating when persistent
  projections, foreground frames, and stable-hull inputs are robust
  under source-side feedback.
\item \textbf{A valid risk or distortion interface.}
  Theorem~\ref{thm:deficit} controls only the supplied score
  \(\mathfrak b_{\mathrm{loc}}\). A sampling model with an estimation-risk
  theorem, or a rate-distortion model with a genuine bottleneck, is
  required before any actual minimax deficit can be transported.
\item \textbf{Iterated perfuming and fixed points.} Once
  \(\mathscr L_\eta\) is applied repeatedly, one faces an entirely new
  dynamical question: do the semigroups converge, cycle, or wander? The
  present paper does not address that long-time behavior.
\item \textbf{Thermodynamic cost of the source-level return.} T-DOME~III
  gives model-dependent reset costs for the calibration procedure at the
  manifest/ego level. Foundation~III has not integrated a specified
  irreversible implementation with the source-side semigroup
  perturbation.
\item \textbf{Multi-source feedback.} Foundation~II handled the
  two-source shared-sector problem at inclusion level, but it did not
  address simultaneous feedback from multiple manifest sectors onto one
  common ambient source dynamics. Cross-source feedback remains open.
\item \textbf{Type~III modular obstacles.} Even if a feedback lift is
  available, one still faces the usual Type~III issues: modular
  invariance of the relevant subalgebras, possible non-uniqueness of
  conditional expectations, central-decomposition effects in
  non-factorial regimes, and the absence of a canonical bounded
  algebra-side generator picture that could be perturbed by finite-rank
  feedback data.
\end{enumerate}

\section{Conditional Input Ledger}\label{sec:ledger}

The theorem chain depends on the following named inputs.

\begin{enumerate}
\item \textbf{FI-A: Foundation~I source-to-manifest interface.}
  The single-source conditional chain from a source datum to a manifest
  pair is available in the sense restated in
  Appendix~\ref{app:interfaces}.
\item \textbf{TD3-A: admissible calibration signal.}
  A one-step natural-gradient update signal \(\delta\sigma\) exists and
  lies inside the T-DOME~III Lyapunov neighborhood.
\item \textbf{FB-A: CP source lift.}
  The calibration signal admits a pointwise \(\sigma\)-weakly continuous
  semigroup lift by normal completely positive unital maps on
  \(\mathcal A_U\).
\item \textbf{FB-B: small-signal differentiability.}
  The lifted feedback semigroup is differentiable at zero coupling so
  that the first-order map \(\Gamma\) is defined; its normalization
  \(\mathfrak G_{\delta\sigma}(\mathbf 1)=0\) is inherited from
  unitality.
\item \textbf{FB-C: chain compatibility.}
  The feedback lift commutes with the ambient source semigroup and
  preserves the named structural inputs needed to rerun the chosen
  Foundation~I chain.
\item \textbf{FB-FD: finite commuting-dephasing witness.}
  Proposition~\ref{prop:fdlift} verifies nontrivial
  \(\delta\sigma\)-dependence, normal UCP dynamics, strict commutation,
  and a source-dynamical change. It is a model check, not a Type~III
  construction.
\item \textbf{DF-A(1): feedback-stable Foundation~I chain.}
  The chosen Foundation~I conditional inputs remain valid for
  \(S_{\mathrm{source}}^{[\eta]}\) on a nontrivial coupling window.
\item \textbf{DF-A(2): manifest Fisher control.}
  The modified manifest-side Fisher sensitivity and ego-rigidity terms
  satisfy the comparison inequalities controlled by
  \(\alpha_F(\eta)\) and \(\beta_{\mathrm{ego}}(\eta)\).
\item \textbf{DF-A(3): supplied local comparison parameters.}
  The coefficient \(\rho(\eta)<\tfrac12\) and the Fisher quantities
  define \(\mathfrak b_{\mathrm{loc}}^{[\eta]}\); no Paper~III
  reconstruction-loss theorem is inherited.
\item \textbf{Downgraded one-step return status.}
  The theorem-level output is a conditionally well-posed one-step return
  map. Its tangent is controlled only at fixed \(t\) as
  \(\eta\downarrow0\), or on windows \(\eta t\ll1\). Iteration,
  convergence, and fixed points remain open.
\end{enumerate}

\appendix
\section{Imported Interfaces}\label{app:interfaces}

This appendix records the external results used by the body in a
self-contained theorem-ready form. The statements are imported
interfaces, not re-proved results.

\begin{remark}[What is and is not imported]
Foundation~III imports a conditional source-to-manifest interface from
Foundation~I and a model-level calibration interface from T-DOME~III.
It does not import the withdrawn Paper~III reconstruction-loss theorem;
only the algebraic comparison formula is retained as supplied data.
Foundation~II is cited in the Introduction solely to locate the later
multi-source feedback problem and does not enter any proof or named
input in the present paper.
\end{remark}

\begin{proposition}[Imported Foundation~I interface]\label{prop:A1}
Let
\[
S_{\mathrm{source}}
=
(\mathcal H_U,\mathcal A_U,\Omega,\{\mathcal E_t\}_{t\ge0})
\]
be a source datum, where
\((\mathcal H_U,\mathcal A_U,\Omega)\) is a Type~III\(_1\) standard-form
von Neumann algebra and
\(\{\mathcal E_t\}_{t\ge0}\) is a pointwise \(\sigma\)-weakly
continuous semigroup of normal completely positive unital maps on
\(\mathcal A_U\). Suppose:
\begin{enumerate}
\item there exists a persistent Nakajima--Zwanzig projection
  \(\mathcal P\) whose projected predual dynamics admits a memory kernel
  \(\mathcal K_*(t,s)\)~\cite{Nakajima1958,Zwanzig1960,BreuerPetruccione2002};
\item a specified finite-coordinate coding model satisfies the
  Foundation~I rank-cost, strict-improvement, and group-action
  assumptions and supplies a compressed foreground frame \(\mathfrak F\);
\item the foreground transfer/stable-hull construction closes so as to
  produce a von Neumann algebra \(\mathcal A_c\subseteq\mathcal A_U\)
  with faithful restricted state
  \(\omega=\langle\Omega,\cdot\,\Omega\rangle|_{\mathcal A_c}\);
\item the admissibility input needed to feed Paper~I~\cite{PaperI} is
  available.
\end{enumerate}
Then the conditional chain
\[
S_{\mathrm{source}}
\to
\mathcal O_A
\to
\mathcal O_B
\to
(\mathcal A_c,\omega)
\]
is defined.
\end{proposition}

\begin{remark}
Foundation~III uses only the existence of the manifest pair produced by
Proposition~\ref{prop:A1}. It does not revisit the internal stable-hull
construction of Foundation~I. The source-datum continuity imported here
is pointwise \(\sigma\)-weak continuity, matching Foundation~I v2. The
published Foundation~I v1 uses the stronger phrase ``strongly
continuous''; the downstream Foundation~I chain used here does not
depend on that stronger condition.
\end{remark}

\begin{proposition}[Conditional T-DOME~III calibration interface]
\label{prop:A2}
Let \(\mathcal O_B\) be a Stage~B output in the sense of
Foundation~I. In the local calibration regime of T-DOME~III there exist:
\begin{enumerate}
\item a natural-gradient update signal
  \(\delta\sigma\in T_\sigma\Sigma\);
\item a Lyapunov neighborhood \(\mathcal U_{\mathrm{cal}}\) on which the
  one-step update stays controlled;
\item a regular prior Fisher term \(\mathcal I_{\mathrm{ego}}\ge0\),
  interpreted as ego rigidity, for the specified local estimation model.
\end{enumerate}
These data determine the one-step manifest-side calibration step used in
the body of the present paper.
\end{proposition}

\begin{remark}
Proposition~\ref{prop:A2} is a one-step interface only. It does not
assert global convergence of repeated calibration updates.
\end{remark}

\begin{definition}[Supplied local Fisher comparison score]
\label{def:A3score}
For stipulated nonnegative quantities \(\mathcal I_F\) and
\(\mathcal I_{\mathrm{ego}}\), and \(\rho\in[0,\tfrac12)\), define
\[
\mathfrak b_{\mathrm{loc}}
:=
\frac{(1-2\rho)\,\mathcal I_F}
{2(\mathcal I_F+\mathcal I_{\mathrm{ego}})}.
\]
The definition makes no assertion that this score lower-bounds an actual
reconstruction loss or estimation risk.
\end{definition}

\begin{remark}[Status of the former Paper~III import]
The original positive drift-floor theorem is withdrawn because its
admissible reconstruction class includes lossless low-dimensional orbit
codes. Definition~\ref{def:A3score} retains only the algebraic expression
needed for Theorem~\ref{thm:deficit}. A future sampling model may support
an estimation-risk bound, and a future bottleneck may support a
rate-distortion bound, but neither result is assumed here.
\end{remark}
 
\newpage
\addcontentsline{toc}{section}{References}

\begin{thebibliography}{99}
\sloppy

\bibitem{FoundationI}
S.~Liu, \emph{Foundation I: From Source to Manifestation --- A Conditional Same-Root Program for Geometry and Modular Flow}, Zenodo (2026), DOI: 10.5281/zenodo.19627786.

\bibitem{FoundationII}
S.~Liu, \emph{Foundation II: Shared Manifestation --- A Conditional Two-Source Program for Shared Accessible Sectors}, Zenodo (2026), DOI: 10.5281/zenodo.19641372.

\bibitem{PaperI}
S.~Liu, \emph{The Inseparability of Geometry and Modular Flow: A Structural Inseparability Theorem and Its Implications for the Hard Problem}, Zenodo (2026), DOI: 10.5281/zenodo.19607387.

\bibitem{PaperII}
S.~Liu, \emph{Higher-Dimensional Consistency of Geometry and Modular Flow: Small-Ball Modular Data and Linearized Gravity}, Zenodo (2026), DOI: 10.5281/zenodo.19607543.

\bibitem{PaperIII}
S.~Liu, \emph{Finite Self-Models and the Open Explanatory Gap: A Conditional Arithmetic Analysis}, Zenodo (2026), DOI: 10.5281/zenodo.19607724.

\bibitem{TDOMEIII}
S.~Liu, \emph{Fisher Information Geometry and Conditional Implementation Costs of Self-Referential Calibration}, Zenodo (2026), DOI: 10.5281/zenodo.18591771.

\bibitem{BreuerPetruccione2002}
H.-P.~Breuer and F.~Petruccione,
\emph{The Theory of Open Quantum Systems},
Oxford University Press (2002).

\bibitem{Nakajima1958}
S.~Nakajima,
\emph{On quantum theory of transport phenomena},
Prog.\ Theor.\ Phys.\ \textbf{20}, 948--959 (1958),
DOI:~10.1143/PTP.20.948.

\bibitem{Zwanzig1960}
R.~Zwanzig,
\emph{Ensemble method in the theory of irreversibility},
J.\ Chem.\ Phys.\ \textbf{33}, 1338--1341 (1960),
DOI:~10.1063/1.1731409.

\bibitem{FrigerioVerri1982}
A.~Frigerio and M.~Verri,
\emph{Long-time asymptotic properties of dynamical semigroups on \(W^*\)-algebras},
Math.\ Z.\ \textbf{180}, 275--286 (1982),
DOI:~10.1007/BF01318911.

\end{thebibliography}
 
\end{document}
